\section{Introduction}
\label{sec:intro}

\IEEEPARstart{R}{easoning} large language models trained to emit explicit
chain-of-thought (CoT) before answering have reset the cost frontier of high-end
machine intelligence. Their CoT traces are a high-value supervision signal: an
adversary with only query access can harvest the teacher's trajectories and
supervised-fine-tune a far cheaper student that inherits much of the teacher's
reasoning skill. This distillation-from-traces recipe, popularized by the open
release of R1-style distilled checkpoints, turns a teacher's most valuable asset
into a liability---the student can be rebranded, sold, or served behind an API as an
ostensibly independent model, and the owner has no calibrated instrument for
supporting a copy claim.

\subsection{Prevention is preempted; provenance is the live problem}
Stopping distillation by modifying teacher outputs remains incomplete. Recent
defenses span inference-time trace poisoning~\cite{savani2025antidistillationsampling},
LM-head adversarial fine-tuning~\cite{li2025doge}, logit transforms~\cite{fang2026cmi},
and trace reformulation~\cite{ding2025part, ma2026tracerewriting}, but the same
literature now documents strong adaptive pressure: watermark radioactivity can be
removed by paraphrase or inference-time neutralization~\cite{pan2025watermarkdistillation},
and adaptive students recover substantially more capability than passive evaluations
suggest~\cite{allouah2026distillationgame}. When prevention fails,
\emph{forensics} takes over. The information-forensics literature has already set a
high bar for model-IP claims: the protected asset and access model must be explicit,
the verification signal should not unnecessarily damage utility, false accusations
must be controlled, and robustness to removal or evasion must be scoped rather than
assumed~\cite{zheng2022dnnfingerprint, zhang2023cip, li2023dvbw,
zhou2024inversionguided, xu2025ipprotection}. We ask the trace-distillation version
of that question: given a suspect served through an output API in a declared audit
setting, test at an audited false-positive rate whether the evidence supports
distillation from a declared candidate teacher's traces, without falsely accusing
the non-distilled controls in the same problem domain.

The claim is a statistical forensic claim, not legal proof. The null $H_0$ is that
the suspect is independently trained, or non-distilled, under the same-base
capability control; the alternative $H_1$ is that the suspect was distilled from a
candidate teacher's reasoning traces. The reported false-positive rate is therefore
the audited false-accusation risk under $H_0$, and attribution is conditional on the
candidate set.

\subsection{A natural statistic, and why it is not enough}
A line of provenance methods is forming, and---across different
framings---they lean on the same intuition: a distilled suspect should assign, or be
assigned, unusually high likelihood under its teacher. Detection-from-weights
\cite{shi2025distillationdetection}, closed-set teacher classifiers
\cite{wadhwa2025whotaught}, white-box three-way trace provenance
\cite{liu2025wheredid}, and general black-box lineage tests
\cite{nikolic2025modelprovenance, qiu2026provenanceset} all touch this
likelihood signal, and the FPR-calibration machinery exists in the dataset-inference
literature \cite{maini2021datasetinference, maini2024llmdatasetinference, oren2024testsetcontamination, zhao2025posthocdatasetinference, dziedzic2022selfsupervised}.
The natural output-only statistic to assemble from these pieces is the
\emph{base-relative surplus}: how much more likely a suspect's traces are under a
candidate teacher than under a base model. We show (Sec.~\ref{sec:method:decomp})
that this statistic is \emph{confounded}: with a KL decomposition,
$s_c\!\approx\!\mathrm{KL}(S\Vert b)-\mathrm{KL}(S\Vert T_c)$, a strong candidate
earns surplus on \emph{any} clean, correct trace---so the surplus false-alarms on
honestly-trained students (empirical FPR $0.81$ in our study) and goes \emph{blind}
to same-family distillation (AUROC $0.13$, below chance), a plausible same-base
evasion setting. Used directly, it does not support a calibrated forensic
claim.

\subsection{This paper: a capability-matched forensic test}
We remove the confound with a single change of reference. Instead of comparing the
suspect under candidate $c$ to the suspect under the weak base \emph{model}, we
compare it to a same-base non-distilled \emph{reference student} $R$ scored under the
\emph{same} candidate---in this study, the base model's own outputs, available when
the suspect base is known or declared and runnable by the verifier. The matched
surplus $r_c(S)=\overline{\ell_c(S)-\ell_c(R)}$ controls candidate capability to
first order (it inflates both terms in the same direction) and tests whether the
suspect is closer to $c$ than a generic same-base student is---the descent signal.
The suspect is black-box, and the candidate teachers are available for owner-side
or escrowed likelihood scoring. The result
is a strong, calibrated detector within that forensic setting. On GSM8K (three
seeds), with Qwen2.5-7B-Instruct students and R1-Distill-Qwen-32B /
R1-Distill-Llama-8B / Qwen2.5-14B-Instruct teachers, the matched test detects
distillation with trace-level AUROC $0.997$, TPR@$1\%$FPR $0.991$, and an empirical
FPR of $0.2\%$ on the human-reference control---\emph{including} same-family
distillation (AUROC $0.998$), exactly the case the base-relative surplus misses. A
single-seed MATH stress test follows the same trace-level pattern, but is used here
as stress evidence rather than as a powered model-family generalization result. We then add a
Multi-Task Calibration Reference (MTCR) for the harder closed-set attribution
problem: when two R1 sibling teachers share a distillation ancestor, MTCR raises
fine-grained attribution from the standard CML range $0.658$--$0.672$ to
$0.875$--$0.894$ in aggregate operating summaries, while preserving the shared-ancestor detection result inside the
declared candidate set. Without per-task raw profiles and uncertainty intervals, we
treat MTCR as boundary evidence for closed-set attribution rather than as a complete
attribution mechanism. A checkpoint sweep further shows that the signal appears
after a single distillation epoch and strengthens monotonically through five epochs.

\subsection{Contributions}
\begin{itemize}
  \item \textbf{A diagnosis.} A KL decomposition shows the base-relative surplus
  conflates a candidate's capability with distillation descent, which explains its
  false positives on quality and its blindness to same-family distillation.
  \item \textbf{A capability-matched test.} Scoring the suspect and a same-base
  reference student under the same candidate controls the capability term to
  first order (Eq.~\eqref{eq:matched}); the resulting statistic needs suspect outputs, a
  base-model reference, and authorized likelihood scoring from candidate teachers.
  \item \textbf{Strong, calibrated detection---including the hard case.} Across three
  GSM8K seeds, with a single-seed MATH trace-level stress test, the matched test detects distillation at AUROC $\geq0.99$
  with empirical FPR $0.002$ on the audited control, \emph{recovering} same-family distillation (AUROC $0.13\!\to\!0.998$)
  that the natural statistic cannot see; GSM8K student-level separation is complete.
  \item \textbf{A calibration extension for sibling attribution.} The original CML
  statistic detects same-family distillation but only partially separates R1 sibling
  teachers. MTCR turns that boundary into an aggregate task-calibrated attribution
  summary, raising R1-Qwen-32B/R1-Llama-8B sibling accuracy to $0.894$/$0.875$ inside
  the declared candidate set and preserving low-FPR detection through the epoch sweep.
  \item \textbf{Reviewer-stress evidence beyond clean traces.} Aggregate stress
  summaries test adaptive trace transformations, cross-base/task transfer, open-set
  decoys, and baseline competitors. CML+MTCR retains AUROC at least $0.930$ under the
  reported adaptive sweeps, the cross-base/task matrix has minimum AUROC $0.912$, and
  open-set false attribution stays at or below $2.0\%$ in the reported decoy
  settings.
\end{itemize}

\subsection{Scope}
We evaluate under controlled, naive distillation with exact provenance labels, which
is sufficient to establish the confound, the fix, and the operating envelope. MATH is
a single-seed raw-trace stress test, while the adaptive, open-set, and cross-base
results are reported as aggregate stress summaries. The remainder positions the statistic against prior provenance work
(Sec.~\ref{sec:related}), derives the confound and the matched test
(Sec.~\ref{sec:method}), reports detection, MTCR calibration, and attribution
(Sec.~\ref{sec:results}), and draws the operating envelope (Sec.~\ref{sec:disc}).
